% !TeX root = kr1.tex
\documentclass[../final.tex]{subfiles}
\begin{document}
% Условия: Downloads/Подготовка_к_кр.pdf, 27.10.26,
% также присланный скриншот codex-clipboard-99032591-db99-414c-b9a9-19332519ffa2.png.
% Решения: Desktop/7.jpg--9.jpg; вычисления проверены и дополнены.
\begin{center}
 {\Large\bfseries Подготовка к контрольной работе}\\[0.5em]
 {\small 27.10.2026}\\[0.6em]
 {\large\bfseries Квантовые газы и тепловое излучение}
\end{center}
\vspace{0.4em}\hrule\vspace{1.2em}
В задачах 1 и 3 используем $k_B=1$. В задаче 2 температура дана
в кельвинах, поэтому постоянную Больцмана сохраняем явно.

\task[1]{Скачок теплоёмкости двумерного бозе-газа}
\condition
Двумерный идеальный бозе-газ находится в потенциале
$U(r)=\beta\sqrt r$, $\beta>0$. Найти скачок теплоёмкости
в точке бозе-конденсации.

\subsection*{Плотность состояний во внешнем поле}
Пусть масса частицы равна $m$, внутреннее вырождение отсутствует.
Одночастичная энергия
$\varepsilon=p^2/(2m)+\beta\sqrt r$.
Число состояний с энергией ниже $\varepsilon$ в квазиклассическом
приближении:
\begin{align*}
 \mathcal G(\varepsilon)
 &=\frac1{(2\pi\hbar)^2}
 \int d^2r\,d^2p\,
 \Theta\left(\varepsilon-\frac{p^2}{2m}-\beta\sqrt r\right)\\
 &=\frac{m}{\hbar^2}
 \int_0^{(\varepsilon/\beta)^2}r(\varepsilon-\beta\sqrt r)\,dr
 =\frac{m\varepsilon^5}{10\hbar^2\beta^4}.
\end{align*}
Здесь при фиксированном $r$ допустимые импульсы занимают круг
площади $2\pi m(\varepsilon-\beta\sqrt r)$.
Дифференцируя, получаем
\[
 D(\varepsilon)=\frac{d\mathcal G}{d\varepsilon}
 =\frac{m\varepsilon^4}{2\hbar^2\beta^4}.
\]
Введём $z=e^{\mu/T}$ и функции Бозе
$g_s(z)=\sum_{l=1}^\infty z^l/l^s$.
С учётом $\int_0^\infty x^{s-1}/(z^{-1}e^x-1)\,dx=\Gamma(s)g_s(z)$:
\[
 N_{\mathrm{ex}}=A T^5g_5(z),\qquad
 E=5A T^6g_6(z),\qquad
 A=\frac{12m}{\hbar^2\beta^4}.
\]
Основной уровень выделен из интеграла; нуль энергии выбран на нём.
В ведущем квазиклассическом приближении его сдвигом относительно
минимума потенциала пренебрегаем.

\subsection*{Температура конденсации и теплоёмкость}
В точке перехода $\mu\to0$, $z\to1$ и $N_{\mathrm{ex}}\to N$:
\[
 N=A T_0^5\zeta(5),\qquad
 T_0=\left[\frac{N\hbar^2\beta^4}{12m\zeta(5)}\right]^{1/5}.
\]
Ниже перехода $z=1$, поэтому
\[
 \frac{N_0}{N}=1-\left(\frac{T}{T_0}\right)^5,\qquad
 C_-=30A\zeta(6)T^5
 =30N\frac{\zeta(6)}{\zeta(5)}\left(\frac{T}{T_0}\right)^5.
\]
Выше перехода из постоянства $N$ и тождества
$zg_s'(z)=g_{s-1}(z)$ находим
\[
 \frac1z\frac{dz}{dT}=-\frac5T\frac{g_5(z)}{g_4(z)}.
\]
Следовательно,
\begin{align*}
 C_+&=\frac{d}{dT}\left[5AT^6g_6(z(T))\right]\\
 &=30AT^5g_6(z)-25AT^5\frac{g_5(z)^2}{g_4(z)}\\
 &=N\left[30\frac{g_6(z)}{g_5(z)}
 -25\frac{g_5(z)}{g_4(z)}\right].
\end{align*}
\answer
Определяем скачок при нагревании как значение сверху минус
значение снизу:
\[
 \boxed{\Delta C=C(T_0+)-C(T_0-)
 =-25N\frac{\zeta(5)}{\zeta(4)}\simeq-23{,}95N}.
\]
Модуль скачка равен $25N\zeta(5)/\zeta(4)$.
В обычных единицах ответ умножается на $k_B$.
Формулы описывают предел большого числа частиц; при конечном $N$
резкого разрыва нет.

\clearpage
\task[2]{Поток фотонов от белого карлика}
\condition
Температура поверхности звезды $T\sim10^4\,\text{К}$,
радиус $R\sim10^4\,\text{км}$. На каком расстоянии плотность
потока фотонов будет $j=1\,\text{см}^{-2}\,\text{с}^{-1}$?
Использовать $\sigma\sim6\cdot10^{-5}
\,\text{эрг}\,\text{см}^{-2}\,\text{с}^{-1}\,\text{К}^{-4}$.

\subsection*{Средняя энергия фотона}
Считаем поверхность абсолютно чёрной, пространство прозрачным,
излучение звезды сферически симметричным.
Планковский спектр даёт среднюю энергию фотона
\[
 \overline\varepsilon_\gamma
 =k_BT\frac{\int_0^\infty x^3/(e^x-1)\,dx}
 {\int_0^\infty x^2/(e^x-1)\,dx}
 =\frac{\pi^4}{30\zeta(3)}k_BT
 \simeq2{,}701\,k_BT.
\]
Для исходящего из чёрной поверхности излучения средняя энергия та же:
геометрический множитель $c/4$ одинаков для потоков энергии и числа
фотонов и сокращается в отношении.

\subsection*{От светимости к числу фотонов}
Поток \emph{энергии} с единицы площади поверхности равен
$\sigma T^4$, поэтому светимость и число испускаемых фотонов в секунду:
\[
 L_\star=4\pi R^2\sigma T^4,\qquad
 \dot N_\gamma=\frac{L_\star}{\overline\varepsilon_\gamma}.
\]
На сфере радиуса $r$ они распределяются по площади $4\pi r^2$:
\[
 j(r)=\frac{\dot N_\gamma}{4\pi r^2}
 =\left(\frac Rr\right)^2
 \frac{\sigma T^3}{2{,}701\,k_B}.
\]
Это поток через площадку, перпендикулярную направлению на звезду.

\subsection*{Численная оценка}
В системе СГС $R=10^9\,\text{см}$,
$k_B=1{,}380649\cdot10^{-16}\,\text{эрг}/\text{К}$. Тогда
\begin{align*}
 r&=R\sqrt{\frac{\sigma T^3}{2{,}701\,k_Bj}}\\
 &\simeq10^9\sqrt{\frac{6\cdot10^{-5}\cdot10^{12}}
 {2{,}701\cdot1{,}380649\cdot10^{-16}\cdot1}}\;\text{см}.
\end{align*}
\answer
\[
 \boxed{r\simeq4\cdot10^{20}\,\text{см}\simeq130\,\text{пк}}.
\]
Так как $r\gg R$, расстояние от поверхности практически равно $r$.

\clearpage
\task[3]{Ферми-газ с дисперсией четвёртой степени}
\condition
Ферми-газ при низкой температуре имеет закон дисперсии
$\varepsilon(p)=\alpha p^4$, $\alpha>0$. Найти $\mu(T)$.

\subsection*{Плотность состояний и энергия Ферми}
Поскольку размерность в условии не указана, используем обычный
трёхмерный однородный газ при фиксированных $N,V$.
Пусть $g$ --- спиновое вырождение ($g=2$ для электронов).
Число состояний в шаре импульсов радиуса $p$:
\[
 \mathcal G(p)=\frac{gV}{(2\pi\hbar)^3}\frac{4\pi p^3}{3}.
\]
Так как $p=(\varepsilon/\alpha)^{1/4}$,
\[
 D(\varepsilon)=B\varepsilon^{-1/4},\qquad
 B=\frac{gV}{8\pi^2\hbar^3\alpha^{3/4}}.
\]
При $T=0$
\[
 N=\frac43B\varepsilon_F^{3/4},\qquad
 p_F=\hbar\left(\frac{6\pi^2n}{g}\right)^{1/3},\qquad
 \varepsilon_F=\alpha p_F^4,\quad n=\frac NV.
\]

\subsection*{Разложение Зоммерфельда при постоянном числе частиц}
При $T\ll\varepsilon_F$ число частиц равно
\begin{align*}
 N&=\int_0^\infty\frac{D(\varepsilon)\,d\varepsilon}
 {e^{(\varepsilon-\mu)/T}+1}\\
 &=\int_0^\mu D(\varepsilon)\,d\varepsilon
 +\frac{\pi^2T^2}{6}D'(\mu)+O(T^4)\\
 &=B\left[\frac43\mu^{3/4}
 -\frac{\pi^2T^2}{24}\mu^{-5/4}
 +O(T^4\mu^{-13/4})\right].
\end{align*}
Полагаем $\mu=\varepsilon_F+\delta\mu$, где
$\delta\mu=O(T^2/\varepsilon_F)$, и вычитаем выражение при $T=0$:
\[
 0=B\varepsilon_F^{-1/4}\delta\mu
 -\frac{\pi^2B T^2}{24}\varepsilon_F^{-5/4}
 +O(BT^4\varepsilon_F^{-13/4}).
\]
\answer
\[
 \boxed{\mu(T)=\varepsilon_F+
 \frac{\pi^2}{24}\frac{T^2}{\varepsilon_F}
 +O\left(\frac{T^4}{\varepsilon_F^3}\right)}.
\]
Для температуры в кельвинах $T^2$ заменяется на $(k_BT)^2$.
Знак поправки положителен: плотность состояний убывает с энергией,
поэтому при нагреве для сохранения $N$ химический потенциал растёт.

\end{document}
